\documentclass[12pt]{article}

\usepackage[utf8]{inputenc}
\usepackage[T1]{fontenc}
\usepackage[brazil]{babel}

\usepackage{geometry}
\geometry{margin=2.5cm}
\usepackage{setspace}
\onehalfspacing

\usepackage{amsmath, amssymb, amsthm}
\usepackage{mathtools}

\numberwithin{equation}{section}

\theoremstyle{plain}
\newtheorem{theorem}{Teorema}[section]
\newtheorem{proposition}[theorem]{Proposição}
\newtheorem{lemma}[theorem]{Lema}
\newtheorem{corollary}[theorem]{Corolário}

\theoremstyle{definition}
\newtheorem{definition}[theorem]{Definição}
\newtheorem{example}[theorem]{Exemplo}
\newtheorem{exercise}{Exercício}[section]

\theoremstyle{remark}
\newtheorem*{remark}{Observação}

\usepackage{hyperref}
\usepackage{csquotes}

\usepackage[
    backend=biber,
    style=authoryear,
    maxcitenames=2
]{biblatex}

\addbibresource{/mnt/c/Users/nicho/Documents/nicholasvoltani.github.io/bibliography.bib}

\newcommand{\R}{\mathbb{R}}
\newcommand{\pref}{\succeq}
\newcommand{\strictpref}{\succ}
\newcommand{\indiff}{\sim}

% ------------------------------------------------
% Title (fill manually)
% ------------------------------------------------
\title{}
\author{}
\date{}

\begin{document}

\hypertarget{conceitos-introdutuxf3rios}{%
\section{Conceitos introdutórios}\label{conceitos-introdutuxf3rios}}

\hypertarget{exercuxedcios}{%
\section{Exercícios}\label{exercuxedcios}}

\hypertarget{ex-1-mas-collell-et-al-exercuxedcio-6.b.4}{%
\subsection{Ex 1 (Mas-Collell et al, exercício
6.B.4)}\label{ex-1-mas-collell-et-al-exercuxedcio-6.b.4}}

O propósito desse exercício é ilustrar como a teoria da utilidade
esperada nos permite tomar decisões consistentes ao lidar com
probabilidades extremamente pequenas ao considerar probabilidades
relativamente grandes. Suponha que uma agência de segurança esteja
pensando em estabelecer um critério sob o qual uma área propensa à
inundação deveria ser evacuada. A probabilidade de inundação é de
\(1\%\). Há quatro resultados possíveis:

\begin{enumerate}
\def\labelenumi{(\Alph{enumi})}
\item
  Nenhuma evacuação é necessária e nada é feito.
\item
  Uma evacuação é feita, que é desnecessária
\item
  Uma evacuação é feita, que é necessária.
\item
  Nenhuma evacuação é feita e uma inundação causa um desastre.
\end{enumerate}

Suponha que uma agência seja indiferente entre o resultado certo de B e
a loteria de A com probabilidade \(p\) e D com probabilidade \(1-p\), e
entre o resultado certo C e a loteria de A com probabilidade \(q\) e D
com probabilidade \(1-q\).\footnote{Tradução: Essa agência possui as
  seguintes preferências: \(B \sim \left(p A + (1-p) D\right)\);
  \(C \sim \left(q B + (1-q) D\right)\); e \(A \succ D\).} Suponha,
também, que ela prefira A a D, e que \(p \in(0, 1)\) e \(q\in(0,1)\).
Assuma que as condições do teorema da utilidade esperada sejam
satisfeitas.

\printbibliography

\end{document}
